\documentclass[11pt]{article}
\usepackage[a4paper,margin=1in]{geometry}
\usepackage{fontspec}
\usepackage[boldfont=false]{xeCJK}
\setCJKmainfont{FandolSong-Regular.otf}
\usepackage{amsmath}
\usepackage{amssymb}
\usepackage{array}
\usepackage{booktabs}
\usepackage{graphicx}
\usepackage{adjustbox}
\usepackage{enumitem}
\setlist[itemize]{topsep=2pt,partopsep=0pt,parsep=0pt,itemsep=2pt,leftmargin=1.4em}
\setlist[enumerate]{topsep=2pt,partopsep=0pt,parsep=0pt,itemsep=2pt,leftmargin=1.6em}
\usepackage{parskip}
\usepackage{fvextra}
\fvset{breaklines=true,breakanywhere=true,fontsize=\small}
\setlength{\parindent}{0pt}

\begin{document}

\begin{center}
  {\LARGE\bfseries The allocation of resources}\\[0.35em]
  {\large Igcse Economics}
\end{center}
\vspace{0.6em}

\section*{Markets and the price mechanism}

Economics is split into two parts:

\begin{itemize}
\item \textbf{microeconomics} 微观经济学 — the study of single parts of the economy, such as one industry or one product.

\item \textbf{macroeconomics} 宏观经济学 — the study of the whole economy, like total jobs, prices, and growth.

\end{itemize}

A \textbf{market} 市场 is any place or system where buyers and sellers meet to trade. It does not need a building — it can be online or over the phone.

When a market is free, no one is in charge of deciding what to make. Instead, the \textbf{price mechanism} 价格机制 (the free movement of prices) \textbf{allocates} 配置 \textbf{resources} 资源 — it decides what is made, how, and for whom. Price does three jobs:

\begin{itemize}
\item the \textbf{signalling} 信号 function — a price change tells buyers and sellers what is happening. A rising price signals "make more"; a falling price signals "make less".

\item the \textbf{incentive} 激励 function — a higher price gives \textbf{firms} 企业 a reason to produce more, because they can earn more.

\item the \textbf{rationing} 配给 function — when a product is scarce, its price rises. The high price shares out the limited amount to those willing to pay.

\end{itemize}

\section*{Demand}

\textbf{Demand} 需求 is the amount of a \textbf{good} 物品 that \textbf{consumers} 消费者 (buyers) are willing and able to buy at each price, in a period of time.

The \textbf{law of demand} 需求定律 says: when the price of a good rises, the quantity demanded falls; when the price falls, the quantity demanded rises. Price and quantity move in opposite directions — an \textbf{inverse relationship} 反向关系. This is true \textbf{ceteris paribus} 其他条件不变 (a Latin phrase meaning "all other things stay the same").

\subsection*{Movements and shifts}

It helps to separate two ideas:

\begin{itemize}
\item \textbf{quantity demanded} 需求量 — the amount bought \emph{at one price}.

\item \textbf{demand} — the whole pattern of buying at every price.

\end{itemize}

On a demand diagram, price is on the up axis and quantity on the across axis. The \textbf{demand curve} 需求曲线 slopes downward from left to right: high price, low quantity; low price, high quantity.

\begin{itemize}
\item A change in the good's \textbf{own price} causes a \textbf{movement along} the demand curve. You slide to a new point on the \emph{same} line.

\item A change in any other condition causes a \textbf{shift} 移动 of the whole curve. A rightward shift means more is demanded at every price; a leftward shift means less.

\end{itemize}

\subsection*{Conditions of demand}

These are the things (other than the good's own price) that shift the demand curve:

\begin{itemize}
\item \textbf{income} 收入 — more income usually means more demand.

\item \textbf{tastes} 偏好 — if a good becomes popular, demand rises.

\item price of \textbf{substitutes} 替代品 — goods used \emph{instead} of each other (tea and coffee). If coffee gets dearer, demand for tea rises.

\item price of \textbf{complements} 互补品 — goods used \emph{together} (cars and fuel). If cars get dearer, demand for fuel falls.

\item \textbf{population} 人口 — more people means more demand.

\item \textbf{advertising} 广告 — good adverts raise demand.

\end{itemize}

\section*{Supply}

\textbf{Supply} 供给 is the amount of a good that firms are willing and able to sell at each price, in a period of time.

The \textbf{law of supply} 供给定律 says: when the price rises, the quantity supplied rises too. Price and quantity move the same way — a \textbf{direct relationship} 正向关系. Firms supply more at a high price because they can earn more.

\subsection*{Movements and shifts}

\begin{itemize}
\item \textbf{quantity supplied} 供给量 — the amount offered \emph{at one price}.

\item The \textbf{supply curve} 供给曲线 slopes upward from left to right.

\item A change in the good's own price → a movement along the curve.

\item A change in any other condition → a shift of the whole curve.

\end{itemize}

\subsection*{Conditions of supply}

\begin{itemize}
\item \textbf{costs of production} 生产成本 — if costs rise, firms supply less (curve shifts left).

\item \textbf{technology} 技术 — better technology lowers costs, so supply rises.

\item \textbf{taxes} and subsidies — a \textbf{tax} 税收 on a good raises firms' costs (supply falls); a \textbf{subsidy} 补贴 (money from the government) lowers costs (supply rises).

\item weather — important for farm goods.

\item number of firms — more firms means more supply.

\end{itemize}

\section*{Price determination}

The price of a good is set where demand and supply meet. Draw both curves on one diagram. They cross at a single point. This point gives:

\begin{itemize}
\item the \textbf{equilibrium} 均衡 price — the price where the amount buyers want to buy exactly equals the amount firms want to sell.

\item the equilibrium quantity — the amount traded at that price.

\end{itemize}

At equilibrium, quantity demanded equals quantity supplied. There is no good left over and no buyer left waiting.

\section*{Price changes}

When a curve shifts, the crossing point moves, so price and quantity change:

\begin{itemize}
\item Demand rises (shifts right): price rises and quantity rises.

\item Demand falls (shifts left): price falls and quantity falls.

\item Supply rises (shifts right): price falls and quantity rises.

\item Supply falls (shifts left): price rises and quantity falls.

\end{itemize}

\subsection*{Disequilibrium}

If the price is \emph{not} at equilibrium, the market is in \textbf{disequilibrium} 非均衡:

\begin{itemize}
\item If price is \textbf{below} equilibrium, buyers want more than firms supply. This is \textbf{excess demand} 超额需求, also called a \textbf{shortage} 短缺. Buyers compete, so the price is pushed up until equilibrium returns.

\item If price is \textbf{above} equilibrium, firms supply more than buyers want. This is \textbf{excess supply} 超额供给, also called a \textbf{surplus} 过剩. Sellers cut the price until equilibrium returns.

\end{itemize}

\section*{Price elasticity of demand (PED)}

\textbf{Price elasticity of demand} 需求价格弹性 measures how much the quantity demanded changes when the price changes.

PED = (\% change in quantity demanded) ÷ (\% change in price)

(We usually ignore the minus sign and read the size of the number.) The \textbf{determinants} 决定因素 (the things that decide PED) include whether the good has close substitutes, whether it is a need or a luxury, and what share of income it takes.

\subsection*{Reading the value}

\begin{itemize}
\item PED greater than 1: demand is \textbf{elastic} 富有弹性. Quantity changes by a large \%. (Goods with many substitutes, like one brand of crisps.)

\item PED less than 1: demand is \textbf{inelastic} 缺乏弹性. Quantity changes by a small \%. (Needs with few substitutes, like petrol or salt.)

\item PED equal to 1: \textbf{unit elastic} 单位弹性. Quantity changes by the same \% as price.

\end{itemize}

\subsection*{PED and total revenue}

\textbf{Total revenue} 总收益 is the money a firm gets from sales: price × quantity sold. (To the buyer, the same money is \textbf{total expenditure} 总支出.) PED tells a firm what happens to revenue if it changes its price:

\begin{itemize}
\item If demand is inelastic, raising the price \emph{raises} total revenue (quantity falls only a little).

\item If demand is elastic, raising the price \emph{lowers} total revenue (quantity falls a lot).

\end{itemize}

This is useful. A firm sets prices knowing how buyers will react. A government taxes inelastic goods (like cigarettes) because people keep buying them, so the tax raises a lot of money.

\section*{Price elasticity of supply (PES)}

\textbf{Price elasticity of supply} 供给价格弹性 measures how much the quantity supplied changes when the price changes.

PES = (\% change in quantity supplied) ÷ (\% change in price)

Supply is elastic if PES is greater than 1, and inelastic if PES is less than 1. The determinants of PES are:

\begin{itemize}
\item \textbf{time} — supply is more elastic over a long time, because firms can change how much they make.

\item \textbf{spare capacity} 闲置产能 — if a firm has unused machines and workers, it can raise output quickly, so supply is elastic.

\item \textbf{stocks} 库存 — goods kept in store can be sold quickly, making supply elastic.

\item ease of switching factors — if a firm can move factors of production to make this good easily, supply is more elastic.

\end{itemize}

Farm goods often have inelastic supply in the short run, because crops take a season to grow.

\section*{The market economic system}

In a \textbf{market economy} 市场经济, the price mechanism makes all the choices. There is no government planning. (The opposite is a \textbf{planned economy} 计划经济, where the government makes the choices.)

\textbf{Advantages} of the market system:

\begin{itemize}
\item Prices guide resources to where they are wanted, without anyone planning.

\item Firms compete, so they cut costs and create new and better goods.

\item Consumers have wide choice.

\end{itemize}

\textbf{Disadvantages}:

\begin{itemize}
\item Some useful goods are not made because they earn no profit.

\item The strong (rich firms and people) do well; the weak may be left with nothing.

\item Markets can fail (see below).

\end{itemize}

\section*{Market failure}

\textbf{Market failure} 市场失灵 happens when the price mechanism leads to an \textbf{inefficient} 低效率 use of resources — resources are not shared out in the best way for society. The main causes:

\subsection*{External costs and benefits}

When you produce or consume something, the cost or benefit to \emph{you} is not always the cost or benefit to \emph{everyone}.

\begin{itemize}
\item a \textbf{private cost} 私人成本 is the cost to the person doing the activity.

\item an \textbf{external cost} 外部成本 is a cost paid by other people — for example, smoke from a factory harms nearby families.

\item the \textbf{social cost} 社会成本 is the private cost plus the external cost (the cost to the whole of society).

\end{itemize}

In the same way:

\begin{itemize}
\item a \textbf{private benefit} 私人效益 is the benefit to the person doing the activity.

\item an \textbf{external benefit} 外部效益 is a benefit that others get for free — for example, your neighbour's tidy garden is nice for you too.

\item the \textbf{social benefit} 社会效益 is the private benefit plus the external benefit.

\end{itemize}

Markets fail because firms and consumers only think about their \emph{private} costs and benefits, not the external ones.

\subsection*{Merit, demerit and public goods}

\begin{itemize}
\item \textbf{merit goods} 有益物品 are better for you than you realise, so the market makes too few (education, healthcare).

\item \textbf{demerit goods} 有害物品 are worse for you than you realise, so the market makes too many (cigarettes, alcohol).

\item \textbf{public goods} 公共物品 are things everyone can use and no one can be stopped from using, like street lights or national defence. The market makes none, because firms cannot charge for them.

\end{itemize}

\subsection*{Abuse of monopoly power}

A \textbf{monopoly} 垄断 is a single firm that controls a market. With no competition, it can charge high prices and make poor-quality goods. This is another market failure.

\section*{The mixed economic system}

A \textbf{mixed economy} 混合经济 has both:

\begin{itemize}
\item a \textbf{private sector} 私营部门 — firms owned by people, guided by the price mechanism.

\item a \textbf{public sector} 公共部门 — activity run by the government.

\end{itemize}

Most real economies are mixed. The government steps in to correct market failure. This is government \textbf{intervention} 干预. The main methods:

\begin{center}
\adjustbox{max width=\linewidth}{%
\begin{tabular}{ll}
\toprule
\textbf{Method} & \textbf{What it does} \\
\midrule
\textbf{maximum price} 最高价格 & a price ceiling, set below equilibrium, to keep a good (like rent) cheap — but it causes a shortage \\
\textbf{minimum price} 最低价格 & a price floor, set above equilibrium, to keep a price up (like a fair wage) — but it causes a surplus \\
\textbf{indirect tax} 间接税 & a tax on a good (added to its price) to cut demand for demerit goods \\
\textbf{subsidy} & money paid to firms to cut the price and raise output of merit goods \\
\textbf{regulation} 监管 & rules and laws, like banning smoking in public \\
\textbf{direct provision} 直接提供 & the government makes and gives out goods itself, like state schools \\
\bottomrule
\end{tabular}}
\end{center}

\subsection*{Judging intervention}

Intervention can fix market failure, but it has costs. Taxes and rules cost money to collect and check. Governments can get things wrong too (this is called government failure). So you should always weigh the good effects against the bad ones before deciding.


\end{document}
